\chaptertitle{第五章\quad 概率初步}

你有没有说过这样的话："天气预报说明天有$80\%$的概率下雨"、"抛硬币正面朝上的概率是$\dfrac12$"、"中彩票头奖的概率几乎为零"。"概率"这个词在生活中无处不在——它衡量的是\textbf{一件事情发生的可能性有多大}。这一章我们从最基础的概率概念开始，逐步建立起计算概率的方法。

\sectiontitle{5.1\quad 确定事件与随机事件}

有些事情不用想就知道结果——"太阳从东方升起"一定会发生，"石头掉到水里会浮起来"一定不会。但有些事的结果是不确定的——"明天会不会下雨"、"抛硬币是正面还是反面"。

\begin{example}
判断下列事件类型：\\
(1) 掷骰子出现$7$点；\\
(2) 今天是星期一，明天是星期二；\\
(3) 从一副扑克牌中随机抽一张，抽到红桃A；\\
(4) 哈尔滨一月的气温低于$0^\circ$C。
\end{example}
\begin{solution}
(1) 骰子只有1~6点——\textbf{不可能事件}。\\
(2) 星期一后面一定是星期二——\textbf{必然事件}。\\
(3) 可能抽到也可能不——\textbf{随机事件}。\\
(4) 哈尔滨一月确实很冷，但暖冬也可能——也属于\textbf{随机事件}（虽然概率非常大）。
\end{solution}

\knowbox{
\textbf{必然事件}：一定会发生的事件，概率为$1$。\\
\textbf{不可能事件}：一定不会发生的事件，概率为$0$。\\
\textbf{随机事件}：可能发生也可能不发生，概率在$0$和$1$之间。}

\begin{exercise}
\item 判断以下事件：\\
(a) 明天太阳从西边升起；\\
(b) 掷一枚硬币，正面朝上；\\
(c) 三角形内角和等于$180^\circ$；\\
(d) 从$1$到$10$中任选一个数，选到$5$。
\item 请各举一个必然事件、不可能事件和随机事件的例子（不能和上面重复）。
\item "一件事发生的概率很小，就说明它不可能发生"——这句话对吗？为什么？
\end{exercise}

\sectiontitle{5.2\quad 古典概型——等可能结果的概率}

最简单的概率问题：抛一枚均匀的硬币，正面朝上的概率是多少？硬币只有两面，均匀硬币两面朝上可能性相同，所以正面朝上的概率是$\dfrac12$。

掷一个均匀的骰子，出现$3$点的概率：$6$个面等可能，所以概率是$\dfrac16$。

这就是\textbf{古典概型}——基本要求是所有可能结果有限，且每个结果等可能。

\knowbox{
\textbf{古典概型}：若随机试验满足：\\
(1) 所有可能结果有限个；(2) 每个结果等可能。\\
则事件$A$的概率：$P(A)=\dfrac{\text{事件}A\text{包含的结果数}}{\text{所有可能结果总数}}$。}

\begin{example}
一个口袋中$2$个红球、$3$个白球（除颜色外完全相同），任取$1$球，求取到红球的概率。
\end{example}
\begin{solution}
所有可能结果：从$5$个球中取$1$个，共$5$种等可能结果。\\
取到红球：口袋中有$2$个红球，所以$2$种结果符合条件。\\
$P(\text{红球})=\dfrac{2}{5}=0.4$。
\end{solution}

\begin{example}
在$1$到$10$中任选一个整数，求选到质数的概率。
\end{example}
\begin{solution}
所有可能结果：$1,2,\dots,10$，共$10$种等可能结果。\\
质数有：$2,3,5,7$（注意$1$不是质数），共$4$个。\\
$P(\text{质数})=\dfrac{4}{10}=0.4$。
\end{solution}

\begin{example}
同时掷两枚均匀硬币，求恰好一枚正面朝上的概率。
\end{example}
\begin{solution}
所有可能结果：正正、正反、反正、反反，共$4$种等可能结果。\\
恰好一枚正面朝上：正反、反正，共$2$种。\\
$P=\dfrac{2}{4}=0.5$。
\end{solution}

\begin{exercise}
\item 一个骰子掷一次，求：(a) 掷出$5$点；(b) 掷出偶数点；(c) 点数大于$4$的概率。
\item 一个盒子中$4$张卡片分别写有$1,2,3,4$，任取一张，求数字大于$2$的概率。
\item 同时掷两枚骰子，$36$种结果中点数之和为$7$的有多少种？概率是多少？
\end{exercise}

\sectiontitle{5.3\quad 树形图与列表法——枚举所有可能}

古典概型的关键是"不重不漏"地列出所有可能结果。当步骤较多时，需要系统的方法——最常用的是\textbf{树形图}和\textbf{列表法}。

\textbf{树形图}适合多步骤试验：

\begin{example}
口袋中$2$个红球（红1、红2）和$1$个白球，先后摸两球（不放回）。用树形图列出所有结果，求两次都摸到红球的概率。
\end{example}
\begin{solution}
第一次有3种可能。第二次只剩2个球。树形图如下：
\begin{center}
\begin{tikzpicture}[scale=0.7]
\node (root) at (0,0){开始};
\node (r1) at (3,2){红1};
\node (r2) at (3,0){红2};
\node (w) at (3,-2){白};
\draw[thick] (root)--(r1); \draw[thick] (root)--(r2); \draw[thick] (root)--(w);
\node (r1r2) at (6,3){红2};
\node (r1w) at (6,1.5){白};
\draw[thick] (r1)--(r1r2); \draw[thick] (r1)--(r1w);
\node (r2r1) at (6,0){红1};
\node (r2w) at (6,-1.5){白};
\draw[thick] (r2)--(r2r1); \draw[thick] (r2)--(r2w);
\node (wr1) at (6,-3){红1};
\node (wr2) at (6,-4.5){红2};
\draw[thick] (w)--(wr1); \draw[thick] (w)--(wr2);
\end{tikzpicture}
\end{center}
所有结果：$(红1,红2)、(红1,白)、(红2,红1)、(红2,白)、(白,红1)、(白,红2)$，共$6$种。\\
两次都红球：$(红1,红2)$和$(红2,红1)$，$2$种。$P=\dfrac{2}{6}=\dfrac13$。
\end{solution}

\textbf{列表法}适合两步试验：

\begin{example}
同时掷两枚骰子，求点数之和为$6$的概率。
\end{example}
\begin{solution}
一枚的结果做行，另一枚做列，交叉处填和：
\begin{center}
\begin{tabular}{c|cccccc}
 & 1 & 2 & 3 & 4 & 5 & 6\\\hline
1 & 2 & 3 & 4 & 5 & 6 & 7\\
2 & 3 & 4 & 5 & 6 & 7 & 8\\
3 & 4 & 5 & 6 & 7 & 8 & 9\\
4 & 5 & 6 & 7 & 8 & 9 & 10\\
5 & 6 & 7 & 8 & 9 & 10 & 11\\
6 & 7 & 8 & 9 & 10 & 11 & 12
\end{tabular}
\end{center}
$36$种结果中，和为$6$的有$5$种：$(1,5),(2,4),(3,3),(4,2),(5,1)$。$P=\dfrac{5}{36}$。
\end{solution}

\begin{center}
\begin{tikzpicture}[scale=0.7]
\node at (2,2.5){\small 和=6的结果（对角线）};
\foreach \x/\y in {1/5,2/4,3/3,4/2,5/1} {
  \filldraw[mainblue] (\x*0.6+0.3,-\y*0.6+3.3) circle (3pt);
}
\draw (0.3,0.3) grid (3.9,3.3);
\node at (0.15,3.0){1}; \node at (0.15,2.4){2}; \node at (0.15,1.8){3};
\node at (0.15,1.2){4}; \node at (0.15,0.6){5}; \node at (0.15,0){6};
\node at (0.6,3.6){1}; \node at (1.2,3.6){2}; \node at (1.8,3.6){3};
\node at (2.4,3.6){4}; \node at (3.0,3.6){5}; \node at (3.6,3.6){6};
\node at (0,4.2){\small 图5-1\quad 两骰子和为6的结果};
\end{tikzpicture}
\end{center}

\knowbox{
\textbf{树形图}：分步展示所有可能，适合$\ge2$步的试验。\\
\textbf{列表法}：用表格列出所有结果，适合两步试验。\\
两种方法的核心目标：不重不漏地枚举。}

\begin{example}
箱子中$3$张卡片（$1,2,3$），先后抽两张（不放回）。用树形图求两张数字之和为偶数的概率。
\end{example}
\begin{solution}
树形图略。所有结果$(1,2),(1,3),(2,1),(2,3),(3,1),(3,2)$共$6$种。\\
和为偶数：$(1,3),(2,2$不存在$),(3,1)$——$(1,3)$和$(3,1)$，共$2$种。\\
$P=\dfrac{2}{6}=\dfrac13$。
\end{solution}

\begin{exercise}
\item 同时掷两枚硬币（$1$元和$5$角），用列表法求：\\
(a) 两枚都正面朝上；(b) 一正一反；(c) 至少一枚正面朝上的概率。
\item 甲、乙玩"石头剪刀布"，假设随机出拳。用列表法求甲获胜的概率。
\item 袋中$2$红$1$白，先后摸两球（不放回），求两次颜色不同的概率。
\end{exercise}

\sectiontitle{5.4\quad 频率与概率——大数定律}

前面讲的古典概型有一个前提：所有结果等可能。但很多实际事件不满足这个条件——比如"明天下雨的概率$80\%$"，是根据历史数据统计出来的。

\begin{example}
历史上数学家们抛硬币的记录：
\end{example}
\begin{center}
\begin{tabular}{c|c|c|c}
实验者 & 抛掷次数 & 正面次数 & 频率\\\hline
蒲丰 & 4040 & 2048 & 0.5069\\
皮尔逊 & 12000 & 6019 & 0.5016\\
皮尔逊 & 24000 & 12012 & 0.5005
\end{tabular}
\end{center}
抛掷次数越多，正面朝上的频率越接近$0.5$。这就是\textbf{大数定律}的直观表现——大量重复试验中，频率稳定在概率附近。

\begin{center}
\begin{tikzpicture}[scale=0.8]
\draw[->,thick] (0,0) -- (7,0) node[right]{抛掷次数};
\draw[->,thick] (0,0) -- (0,3.5) node[above]{频率};
\foreach \y in {0.48,0.5,0.52} {
  \draw (0.1,(\y-0.46)*100) -- (-0.1,(\y-0.46)*100) node[left]{$\y$};
}
\draw[dashed] (0,(0.5-0.46)*100) -- (7,(0.5-0.46)*100);
\node at (6.3,(0.5-0.46)*100+0.2){$0.5$};
\draw[mainblue,thick] plot coordinates {
  (0.5,0.6) (1,1.2) (1.5,0.8) (2,2.5) (2.5,1.8) (3,2.2) (3.5,1.5) (4,2) (4.5,1.6) (5,1.8) (5.5,1.7) (6,2) (6.5,1.6)
};
\node at (3.5,-0.8){\small 图5-2\quad 频率随抛掷次数的变化（示意图）};
\end{tikzpicture}
\end{center}

\knowbox{
\textbf{频率与概率的关系}：\\
大量重复试验中，事件发生的频率会稳定在某个常数附近——这个常数就是概率。\\
试验次数越多，频率越接近概率。这就是\textbf{大数定律}。}

\begin{example}
小明掷硬币$100$次，正面$54$次。频率是多少？掷$1000$次会怎样？
\end{example}
\begin{solution}
频率$=54\div100=0.54$。掷$1000$次时，频率通常会比$0.54$更接近$0.5$——因为次数越多频率越稳定。
\end{solution}

\begin{example}
商场抽奖：$1000$张奖券中一等奖$1$张、二等奖$10$张、三等奖$50$张。求中奖和中一等奖的概率。
\end{example}
\begin{solution}
中奖券$=1+10+50=61$张，$P(\text{中奖})=\dfrac{61}{1000}=0.061$。\\
$P(\text{一等奖})=\dfrac{1}{1000}=0.001$。
\end{solution}

\begin{exercise}
\item 某射手射击$100$次命中$85$次，频率是多少？如果射击$10000$次，你估计命中率会稳定在多少附近？
\item 天气预报"明天下雨概率$80\%$"，这句话的意思是：\\
A. 明天$80\%$的时间下雨；B. 明天一定会下雨；C. 类似气象条件下约$80\%$的日子会下雨。
\item 袋中有$10$个球（红+白，数量未知）。小明摸了$50$次（放回），红球$32$次。估计红球有几个？
\end{exercise}

\sectiontitle{5.5\quad 期望值——概率的平均}

概率除了描述"可能性"，还能用来计算"平均结果"。这就是\textbf{期望值}。

\begin{example}
掷一枚骰子，点数的期望值是多少？
\end{example}
\begin{solution}
每个点数$1,2,3,4,5,6$出现的概率都是$\dfrac16$。\\
期望值$=1\times\dfrac16+2\times\dfrac16+3\times\dfrac16+4\times\dfrac16+5\times\dfrac16+6\times\dfrac16
=(1+2+3+4+5+6)\times\dfrac16=3.5$。\\
期望值$3.5$虽然不是实际可能掷出的点数，但如果掷很多次，平均点数就会接近$3.5$。
\end{solution}

\begin{example}
抽奖：$100$张奖券中，一等奖$1$名奖$500$元，二等奖$5$名各奖$20$元，其余无奖。每张奖券$10$元。从期望值角度看，买一张划算吗？
\end{example}
\begin{solution}
中一等奖概率$=1/100$，中二等奖$=5/100$，不中奖$=94/100$。\\
期望收益$=500\times\dfrac1{100}+20\times\dfrac5{100}+0\times\dfrac{94}{100}=5+1+0=6$元。\\
买一张花$10$元，期望回报只有$6$元——所以从期望值来看不划算。这就是彩票的原理：多数人赔，少数人赚。
\end{solution}

\knowbox{
\textbf{期望值}：随机变量各可能值乘以其概率后的总和。\\
$E(X)=x_1p_1+x_2p_2+\cdots+x_np_n$。\\
可以理解为"大量重复试验中的平均结果"。}

\begin{exercise}
\item 抛一枚硬币，约定正面赢$2$元，反面输$1$元。求期望收益。这个游戏公平吗？
\item 同时掷两枚骰子，点数之和的期望值是多少？（提示：可以利用对称性）
\item 某保险产品：一年保费$500$元，若出险赔付$10$万元。已知出险概率$0.3\%$。从期望值看，保险公司平均从每个客户身上赚多少钱？
\end{exercise}

\sectiontitle{本章小结}

概率衡量事件发生的可能性。必然事件概率$1$，不可能事件概率$0$，随机事件概率在$0$到$1$之间。古典概型是最基本的概率模型：$P(A)=\dfrac{\text{有利结果数}}{\text{总结果数}}$。复杂试验用树形图或列表法枚举所有可能。大量重复试验中，频率稳定于概率——这是大数定律。期望值则告诉我们"长期平均结果"是多少，是概率用于决策的重要工具。

\sectiontitle{课后练习}

\subsection*{A组\quad 基础巩固}

\begin{exercise}
\item 判断必然事件、不可能事件还是随机事件：\\
(a) 掷骰子出现$1$点；(b) 没有水分种子发芽；(c) 买彩票中一等奖；(d) 标准大气压下$100^\circ$C水沸腾。
\item $3$个红球$5$个白球，任意摸$1$个，求摸到白球的概率。
\item 从$1,2,3,4,5$中任取一个数，求它既是奇数又是质数的概率。
\item 同时掷两枚硬币，求两枚都反面朝上的概率。
\item 一个骰子掷一次，求点数小于$3$的概率。
\item 某射手射击$50$次命中$42$次，命中频率是多少？
\end{exercise}

\subsection*{B组\quad 挑战提升}

\begin{exercise}
\item 一个家庭有两个孩子（生男生女概率相同）。用树形图列出所有性别组合，求性别不同的概率。
\item $4$张卡片$A,B,C,D$，一次抽两张，求恰好抽到$A$和$B$的概率。
\item 甲、乙各掷一枚骰子比大小。甲掷出$5$点的概率？甲点数大于乙的概率？（用列表法）
\item 小明抛硬币$100$次得$48$次正面。小红说"次数太少，我抛$1000$次"得$495$次。谁的频率更接近$0.5$？说明什么？
\item 密码$4$位数字（$0$~$9$），随机输入一次猜对的概率？如果是$4$位字母（$A$~$Z$）呢？
\item 设计一个游戏：同时掷两枚骰子，规定什么和算甲赢、什么和算乙赢——使得双方获胜概率相等。
\item 小明和小红玩：同时抛两枚硬币，都正面小明得$3$分，都反面小明得$1$分，一正一反小红得$2$分。小明说"我有两种方式得分，你只有一种，所以我赢面大"。这个游戏公平吗？用期望值说明。
\end{exercise}

\newpage
